% SPDX-License-Identifier: GPL-3.0-or-later
% Copyright (C) 2026 Alessandro Greco
\section*{Глава \ref{cap:15-dos}: защита от DoS}

\begin{soluzione}{es:15-bucket}
В первую минуту полный резервуар даёт 60 токенов, а пополнение добавляет $2 \times 60 = 120$, то есть в сумме примерно 180 цепей; фактически несколько меньше, потому что клиент помечается, как только резервуар опустеет. За час при установившейся скорости получится примерно $60 + 2 \times 3600 = 7260$ цепей. Обычный клиент создаёт лишь несколько десятков цепей.
\end{soluzione}

\begin{soluzione}{es:15-guard}
Только первый ретранслятор видит настоящий IP-адрес клиента: на следующих переходах \enquote{адрес} канала принадлежит предыдущему ретранслятору. Если промежуточный ретранслятор ограничивает число цепей для каждого адреса, он фактически ограничивает все цепи входного узла, который обслуживает тысячи пользователей, и наносит ущерб всем, а не атакующему. Поэтому ограничения применяются только к неаутентифицированным соединениям клиентов.
\end{soluzione}

\begin{soluzione}{es:15-amplificazione}
Клиент отправляет сообщение \texttt{INTRODUCE1} по цепи, которую может использовать повторно. Сервис должен расшифровать сообщение, выполнить рукопожатие hs-ntor и построить новую цепь как минимум из трёх переходов к точке встречи. Это требует соответствующих рукопожатий ntor с участвующими ретрансляторами, а также ячеек создания и расширения для каждого перехода. Одна ячейка клиента тем самым вызывает несколько операций с открытым ключом и десятки ячеек в сети: усиление нагрузки на один или два порядка.
\end{soluzione}

\begin{soluzione}{es:15-bit}
При требовании $k$ нулевых бит можно задавать только уровни трудоёмкости, равные степеням двойки: 1, 2, 4 и так далее. При переходе на следующий уровень вычислительная стоимость удваивается. Линейная проверка позволяет задать любое целое значение трудоёмкости $E$. Регулятор с обратной связью, изменяющий трудоёмкость малыми шагами, нуждается в более точной шкале: при удвоениях он колебался бы между слишком низкими и слишком высокими значениями.
\end{soluzione}

\begin{soluzione}{es:15-sforzo}
Нет. Трудоёмкость $E$ входит в задачу, по которой вычисляются и решение Equi-X, и хеш $R$. При объявлении $E = 5000$ задача меняется, поэтому решение, найденное для $E = 50$, больше недействительно. Клиент должен был выбрать $E = 5000$ до начала поиска.
\end{soluzione}

\begin{soluzione}{es:15-oom}
Лабораторное упражнение. Ячейки накапливаются в очереди цепи промежуточного ретранслятора, пока общий объём не превысит \opt{MaxMemInQueues}. Тогда диспетчер памяти закрывает цепи с наиболее старыми данными в очередях, в том числе нашу, и записывает в журнал объём освобождённой памяти.
\end{soluzione}

\begin{soluzione}{es:15-coda}
Упражнение по программированию. Ожидается, что запросы с высокой трудоёмкостью будут иметь короткое и устойчивое время ожидания, а запросы с нулевой трудоёмкостью будут обслуживаться только при почти пустой очереди или отбрасываться при её разгрузке. Это намеренный результат: он защищает пользователей, готовых затратить немного вычислительных ресурсов.
\end{soluzione}

\section*{Глава \ref{cap:16-censura}: цензура, мосты и подключаемые транспорты}

\begin{soluzione}{es:16-guard}
Мост (bridge) выполняет роль входного узла (guard), первого перехода цепи, но не фигурирует в публичном консенсусе. Клиент узнаёт его адрес из полученной частным образом строки конфигурации моста. При \opt{UseBridges} клиент использует мосты \emph{вместо} публичных входных узлов: добавление обычного входного узла означало бы прямое подключение к опубликованному адресу, то есть именно то, что цензор блокирует или отслеживает.
\end{soluzione}

\begin{soluzione}{es:16-strategie}
\enquote{Не быть похожим ни на один известный протокол}: obfs4 (для блокирования цензору пришлось бы запретить весь нераспознаваемый трафик). \enquote{Быть похожим на разрешённый трафик}: meek (потребовалось бы заблокировать крупную сеть доставки содержимого), Snowflake (заблокировать WebRTC, брокер и множество добровольных прокси), WebTunnel (блокировать HTTPS к неизвестным сайтам либо сайтам, на которых размещены мосты), Conjure (блокировать широкие диапазоны адресов у сотрудничающих провайдеров).
\end{soluzione}

\begin{soluzione}{es:16-sondaggio}
Цензор может искать долгие подключения к адресам без связанного доменного имени, высокую энтропию байтов уже в первом пакете, характерные распределения размеров и интервалов передачи (например, всплески, связанные с ячейками Tor), а также отношение входящего и исходящего трафика, напоминающее просмотр веб-страниц. \texttt{iat-mode} добавляет случайные задержки и фрагментацию, затрудняя распознавание размеров и времени ценой снижения производительности.
\end{soluzione}

\begin{soluzione}{es:16-lox}
Цензор блокирует полученные мосты, и заблокированный мост не позволяет продвинуться на следующий уровень: соответствующие учётные данные остаются на низком уровне и не дают права приглашать других. Цензор может не блокировать найденные мосты, пока не достигнет высокого уровня, а затем заблокировать их все сразу. Однако для этого пришлось бы месяцами поддерживать множество терпеливых идентичностей, а мосты всё это время обслуживали бы обычных пользователей. Кроме того, каждые учётные данные высокого уровня дают доступ лишь к небольшому числу мостов.
\end{soluzione}

\begin{soluzione}{es:16-entropia}
Оба захвата имеют энтропию, близкую к восьми битам на байт: TLS после рукопожатия тоже передаёт зашифрованные данные. Различия заключаются в нескольких байтах: заголовках записей TLS (\texttt{16 03 01}, \texttt{17 03 03}, затем длина) и первоначальном незашифрованном рукопожатии. Система глубокой проверки пакетов (DPI) распознаёт эти признаки без статистического анализа.
\end{soluzione}

\begin{soluzione}{es:16-sondaggio-lab}
Лаборатория~1 завершается по времени ожидания (\texttt{TimeoutError}), потому что мост не отвечает. При неверном параметре \texttt{cert} клиентская программа \cmd{obfs4proxy} не может завершить рукопожатие; мост не распознаёт метку аутентификации клиента, не отвечает и через некоторое время закрывает соединение. В журнале моста могут появиться только сообщения о неудачном рукопожатии, в зависимости от уровня журналирования.
\end{soluzione}

\section*{Глава \ref{cap:17-torrc}: установка и torrc}

\begin{soluzione}{es:17-verifica}
tor сообщает о неизвестной опции сообщением вида \texttt{Unknown option '...'. Failing.} и завершает работу, не запускаясь. Команда \texttt{--dump-config short} для клиентского файла \file{torrc} показывает лишь параметры, отличающиеся от значений по умолчанию: в лабораторной сети это, среди прочего, \opt{TestingTorNetwork}, записи \opt{DirAuthority} и порты.
\end{soluzione}

\begin{soluzione}{es:17-isolamento}
Если почта и форум используют одну цепь, выходной ретранслятор видит примерно в одно время по этой цепи соединение с почтовым сервером пользователя (часто связанное с его именем) и с форумом под псевдонимом. Он может связать псевдоним с человеком. Решения: два разных SOCKS-порта, по одному для каждого приложения (разные порты изолированы по умолчанию, если им не назначена общая \opt{SessionGroup}); либо разные учётные данные SOCKS с \texttt{IsolateSOCKSAuth}; либо \texttt{IsolateDestAddr} на общем порту.
\end{soluzione}

\begin{soluzione}{es:17-exitnodes}
Атакующий, наблюдающий сети Италии, мог бы увидеть и пользователя, входящего в сеть Tor, и итальянский выходной ретранслятор, подключающийся к итальянскому сайту. Оба конца находятся в одной наблюдаемой области — это как раз сценарий корреляции, от которого модель угроз Tor не гарантирует защиты. Кроме того, итальянских выходных ретрансляторов мало; ограничение выбора делает пользователя более узнаваемым и облегчает атаки со стороны злонамеренных ретрансляторов.
\end{soluzione}

\section*{Глава \ref{cap:18-stem}: порт управления и Stem}

\begin{soluzione}{es:18-codici}
\texttt{552} — постоянная ошибка в ответ на команду: ключ \texttt{pippo} не существует. \texttt{650} — не ответ на команду, а асинхронное событие, передаваемое после подписки клиента с помощью \texttt{SETEVENTS}. Оно может поступить в любой момент, в том числе во время ожидания ответа на другую команду.
\end{soluzione}

\begin{soluzione}{es:18-onion}
План: вызвать метод контроллера \file{create\_ephemeral\_hidden\_service}, сопоставив виртуальный порт 80 с локальным портом 8080 и указав \file{await\_publication}; метод возвращает адрес onion-сервиса. Подключиться к нему можно командой \texttt{curl --socks5-hostname 127.0.0.1:9008 http://<address>.onion/}. Ключи временного сервиса tor хранит в памяти и возвращает контроллеру, который может сохранить их для последующего восстановления того же адреса. При закрытии управляющего соединения сервис исчезает, если не задать режим независимой работы (detach).
\end{soluzione}

\begin{soluzione}{es:18-banda}
План: зарегистрировать обработчик для \texttt{EventType.BW}, получающий каждую секунду количество прочитанных и записанных байтов (\texttt{event.read} и \texttt{event.written}), и выводить значения во время загрузки. Скорость приёма растёт до предела, определяемого механизмом управления перегрузкой и ресурсами компьютера.
\end{soluzione}

\begin{soluzione}{es:18-percorso}
При \file{\_\_LeaveStreamsUnattached} tor не прикрепляет потоки к цепям автоматически. Команда \texttt{EXTENDCIRCUIT 0 test004r,test005r,test006e} строит цепь и возвращает её идентификатор. Когда поступает событие \texttt{STREAM} со статусом \texttt{NEW}, вызывают \texttt{ATTACHSTREAM}. Ручной выбор маршрута опасен: он обходит взвешенный выбор ретрансляторов и правила использования входных узлов, делает пользователя узнаваемым и может позволить вредоносной программе направить его трафик через подконтрольные противнику ретрансляторы.
\end{soluzione}
